\documentclass[10pt]{article}
\usepackage[utf8]{inputenc}
\usepackage{MyPreambles}
\usepackage{import}
\usepackage{hyperref}
\usepackage{multicol}
\usepackage{paracol}


\setlength{\headheight}{35pt}
\setlist{noitemsep}
\setlength\intextsep{0pt}
\setlist[description]{noitemsep, topsep=0pt, leftmargin=*}


\usepackage[dvipsnames]{xcolor}
\usepackage{fontspec} % if using XeLaTeX/LuaLaTeX
% \usepackage{inconsolata} % for nicer \ttfamily with pdfLaTeX
\usepackage{minted}

\definecolor{codebg}{RGB}{245,245,248}


\title{Notes on DS 4400: Machine Learning}
\author{William Wei}
\date{September 2026}

\begin{document}

\maketitle

\tableofcontents

\pagebreak


 
\pagebreak
\section{Unit 1: MGF and Others }
\subsection{Topic 1: Moment Generating Functions}
\begin{multicols}{2}
\begin{definition}[Moment]
The $k$th moment of a random variable $X$ is $E(X^k)$. It only exists if it is finite. 
\end{definition}
Note that the first moment is just the expected value of some RV. 
\begin{definition}[MGF]
The moment generating function (MGF) of a random variable $X$ is 
$$
M_X(t) = E(e^{tX}),
$$
provided that the expectation is finite for every $t$ in some open neighborhood of $0$. 
\end{definition}
The function generates moments of the RV because of the following theorem. 
\begin{theorem}
The $k$th moment of a random variable $X$ is the $k$th derivative of its MGF evaluated at $t = 0$. That is, 
$$
E(X^k) = \left[M_X(t)^{(k)}\right]_{t = 0}
$$
\end{theorem}
\begin{proof}
Recall that for a function $f(x)$, its Maclaurin series is 
$$
f(x) = \sum_{k = 1}^\infty \frac{x^k}{k!}f'(0).
$$
We expand the MGF and take its derivative. 
\begin{align*}
    M_X^{(k)}(t) &= \left[E(e^{tX})\right]^{(k)} \\
    &= \left[E\left(\sum_{i = 0}^\infty \frac{(tX)^i}{i!}\right)\right]^{(k)} \\
    &= \left[\sum_{i = 0}^\infty E\left(X^i\right)\frac{t^i}{i!}\right]^{(k)}
\end{align*}
At $t = 0$, all terms containing a non-zero power of $t$ vanishes, leaving behind only the term at $i = k$, which is precisely $E(X^k)$. 
\end{proof}
\begin{remark}
There are a lot of tricks of finding the MGF of a RV. We list a few common ones below: 
\begin{itemize}
    \item Using L'Hopital's rule to check limit at $0$. 
    \item Recognizing that a series or integral is just a different from of a common series or integral. 
    \item Realizing that the PDF should sum/integrate to $1$ over the sample space. 
\end{itemize}
    
\end{remark}

\begin{theorem}
Some properties of MGF
\begin{itemize}
    \item $M_x(t) = M_Y(t) \implies X = Y$
\item $M_{aX+b}(t) = e^{bt}M_X(at)$
    \item $M_{X+Y} = M_X(t) M_Y(t)$
\end{itemize}
\end{theorem}
Below we present a table of commonly use MGF's. 
\begin{tabular}{c|c|c|c}
    Distr& Param & MGF & Domain \\
    \hline\hline
    Bernoulli & $p$ & $1 - p + pe^t$ & all $t$ \\
    Binomial & $n, p$ & $(1 - p + pe^t)^n$ & all $t$ \\
    Poisson & $\lambda$ & $\exp{[\lambda(e^t - 1)]}$ & all $t$ \\
    Normal & $\mu, \sigma^2$ & $\exp[\mu t + \sigma^2t^2/2]$ & all $t$\\
    Exp & rate $\lambda$ & $\lambda / (\lambda - t)$ & $t < \lambda$ \\
    $\chi^2$ & $\nu$ & $(1 - 2t)^{-\nu/2}$ & $t < 1/2$
\end{tabular}
\end{multicols}
%\pagebreak
\section{Dynamic Programming and Greedy Methods}

\subsection{Lecture 6: Dynamic Programming 1}
\begin{multicols}{2}
Dynamic programming (DP) only differ from divide and conquer (DaC) algorithms in that \textbf{DaC divide up problems into separate subproblems, while DP divide up problems into overlapping subproblems.}

If we have already calculated a value before, we should not forget it. We store those values into a table and refer to it later when needed. 


\begin{leftbar}
\textbf{Weighted Interval Scheduling Problem}
\begin{itemize}
        \item Give a list of $n$ jobs $J_i$
        \item $J_i$ has a start $s_i$ and finish time $f_i$.
        \item If two jobs' time rages overlap, they are incompatible. 
        \item Each job comes with an associated weight $w_i$, representing their value.
        \item \textbf{Goal} Find the \textit{maximum weight subset} of mutually compatible jobs.
    \end{itemize}
\end{leftbar}
\begin{analysis}
    We need to figure out how to break it up into subproblems. Brute force gives $2^n$ subsets, which is huge. If we can solve for the first $i -1$ jobs, perhaps we can solve until job $i$. Therefore, \textit{we sort the jobs by their finish time $f_i$}. Define $OPT(i)$ as the \textit{maximum weight subset possible} using jobs $1$ to $i$. Our new \textbf{goal} is to find $OPT(i)$ in terms of $OPT(1), \ldots, OPT(i - 1)$. 

    To find $OPT(i)$, 
    
    \begin{itemize}
        \item If we do not want to pick $J_i$, then $OPT(i) = OPT(i - 1)$. 
        \item If we do wish to pick $J_i$, then we need to pick the maximum jobs from $1 ,\ldots, i - 1$ that do not conflict with $J_i$. We find the largest index $j < i$ that does not conflict with $i$. Redefine this index $j = p_i$, which can be found via binary search. 
    \end{itemize}
    This process gives a Bellman Optimality Equation 
    $$
    OPT(i) = 
    \begin{cases}
        0 & i = 0 \\
        \max(OPT(i - 1), w_i + OPT(p_i)) & i > 0
    \end{cases}
    $$
\end{analysis}
This will take a total of $O(n \log n)$ time. 

\begin{leftbar}
    \textbf{Maximum Grid Path Sum}
    \begin{itemize}
        \item \textbf{Input} an $n \times m$ grid of positive numbers.
        \item The path starts on top left, only moves right or down on each step, ends on bottom right.
        \item We sum the elements along the path
        \item \textbf{Goal} Find maximum path sum among all sums.
    \end{itemize}
\end{leftbar}

\begin{analysis}
    Brute force gives about $2^{2n}$ paths for an $n \times m$ grid, which is too large. We use the same idea as we did in WIS, that is, dividing the problem into two subproblems, namely $(n -1) \times m$ and $n \times (m - 1)$. If the grid has size $1 \times 1$, the solution is trivial. If the size was $2\times 2$, there are only 2 solutions. Then when the sizes become $2\times 3$ or $3\times2$, there are only 3 possible solutions. We therefore try to reduce the problem into smaller grids. Below is the Bellman Optimality Equation
    $$
    S(n,m) = A[n,m] + \max(S(n - 1, m), S(n, m - 1)),
    $$
    where $A[n,m]$ is the $(n,m)$ element, which is on the bottom right corner. The intuition is that, if we were to come to the bottom right corner from the left, we would use the value $S(n - 1, m)$. If we came from the top, we would use the value $S(n, m - 1)$. For a visual, we can populate the grid with the maximum sum path that will lead to that point. Just apply the $S(n,m)$ formula to the desired point $(i, j)$. 

    Note that this algorithm using DP will only take $O(mn)$ time and $O(n)$ space. 
\end{analysis}
\end{multicols}

\newpage
\subsection{Lecture 7: Dynamic Programming 2 - Sequences}
\begin{multicols}{2}
\begin{definition}[Sequences]
    A sequence $s_n$ is a function from the set $\{n \in \Z : n \geq m\}$ to anything. A sequence $t_n$ is a subsequence of $s_n$ if and only if there exists an increasing map $\sigma : \N \to \N$ such that $\sigma(k) = n_k$ for $k \in \N$. That is, $t_k = t(k) = s\circ \sigma(k) = s(n_k) = s_{n_k}$. 
\end{definition}

\begin{leftbar}
    \textbf{Longest Increasing Subsequence (LIS)}
    
    Given a sequence of integers, what is the length of the longest increasing subsequence?
\end{leftbar}
\begin{analysis}
    We apply the strategy of dynamic programming again. For the input $x = (x_1, \ldots, x_n)$, define $LIS(j)$ to be the length of the LIS of $x_1, \ldots, x_j$. Note that for $LIS(j+1)$, if we naively compare $x_j$ and $x_{j+1}$, it does not solve the problem, since the increasing \textit{subsequence} may not include $x_j$. Therefore, we enforce that $LIS(j)$ must end with $x_j$. This means that the new LIS might be shorter than the real LIS at $j$, but this convention enables us to recurse up to $j$ and obtain our results. 

    Suppose there is a $i < j$ such that $x_i < x_j$, then we append $x_j$ to the end of the LIS ending with $x_i$. This appending operation will increase the length by $1$. If we perform this operation for all $i < j$ and $x_i < x_j$, we have succeeded in recusing. Mathematically, 
    $$
    LIS(j) = 
    \begin{cases}
        1 & j = 0 \\
        1 + \max_{1 \leq i < j, x_i < x_j}LIS(i) & \textrm{otherwise}
    \end{cases}
    $$
    The notation could be confusing. $LIS(n)$ is not what we want for a sequence of length $n$. This is because we enforced that $LIS(n)$ to be the length of the increasing subsequence that end with $x_n$, which may or may not be included in the actual longest increasing subsequence. What we truly need is $\max_{1 \leq i \leq n} LIS(i)$. 

    The runtime of this algorithm is $O(n^2)$. This is because we need $O(i)$ time to find $LIS(i)$, and $\sum_{1\leq i \leq n}O(i) = O(n^2)$. Same idea from this algorithm also works for decreasing, non-increasing, and non-decreasing subsequences. 
\end{analysis}
\begin{leftbar}
    \textbf{Sequence Similarity Problem}
    
    Given two sequences, how do we measure how close they are? 
\end{leftbar}
\begin{definition}[Edit Distance]
    Given two strings $x\in \Sigma^n$ and $y \in \Sigma^m$, the edit distance is the number of insertions, deletions, and substitutions required to turn $x$ into $y$. 
\end{definition}
\begin{definition}[Alignment and cost]
    An alignment is a mapping from $x$ to the a subsequence of $y$. Given an alignment between two strings, the cost of such alignment is the number of mismatched positions. 
\end{definition}
Note that the edit distance can also be defined as the cost of minimum cost alignment. 

\begin{analysis}
    In an alignment of $x$ and $y$, the last column could have the following possibilities. 
    \begin{itemize}
        \item $(x_i, \cdot)$: delete from $x$
        \item $(\cdot, y_i)$: insert into $x$
        \item $(x_i, y_i)$: matching characters
        \item $(x_i, y_i)$: mismatching characters
    \end{itemize}

    Define $OPT(i, j)$ to be the edit distance of strings $x = x_1\cdots x_i$ and $y = y_1\cdots y_j$. We analyze the last column in cases again: 
        \begin{itemize}
        \item $(x_i, \cdot)$ costs $1 + OPT(i - 1, j)$. 
        \item $(\cdot, y_j)$ costs $1 + OPT(i, j - 1)$.
        \item matching $(x_i, y_j)$ costs $OPT(i - 1, j - 1)$.
        \item mismatching $(x_i, y_j)$ costs $1 + OPT(i - 1, j - 1)$.
    \end{itemize}
    It is natural to take the minimum of all of those four cases above, because we want the minimum amount of edits to transform $x$ into $y$. Sometimes deletion could benefit us in the short term, but it could actually just be a mismatching character. This is why we take the minimum of all of them, so that we try out all of them and find out the least number of edits. Define as a shorthand 
    $$
    \min_{i, j} opt(i,j) = \min
    \begin{cases}
        1 + OPT(i - 1, j) \\
        1 + OPT(i, j -1) \\
        OPT(i - 1, j -1)
    \end{cases} 
    $$
    Note that if $(x_i, y_j)$ mismatch, we must add an operation outside of the $\min$, since no matter what the minimum edits the previous OPTs give us, we still need to make an operation in the current step. Define again a shorthand that 
    $$
    1 [P] = 
    \begin{cases}
        1 & P \\
        0 & \textrm{not } P
    \end{cases}
    $$
    As a general recursion formula, we have 
    $$
    OPT(i, j) = 
    \begin{cases}
        i & j = 0 \\
        j & i = 0 \\
        1[x_i\neq y_j] + \min_{i,j}opt(i,j) & \textrm{otherwise}
    \end{cases}
    $$
\end{analysis}
\begin{example}
    The edit distance between \say{SALT} and \say{SMALL} is $2$. This can be seen in the table below 

    \begin{center}
        
    \begin{tabular}{c|c c c c c c}
         &  & S & M & A & L & L \\
         \hline 
         & 0 & 1 & 2 & 3 & 4 & 5 \\
        S & 1 & \textbf{0} & \textbf{1} & 2 & 3 & 4 \\
        A & 2 & 1 & 1 & \textbf{1} & 2 & 3 \\
        L & 3 & 2 & 2 & 2 & \textbf{1} & 2 \\
        T & 4 & 3 & 3 & 3 & 2 & \textbf{2}
    \end{tabular}
    \end{center}

    
    Note that we can also backtrack to see where we need to make modifications to \say{SMALL} by looking at the path that led us to the number $2$ on the bottom right corner. 
\end{example}
\end{multicols}

\subsection{Lecture 8: Dynamic Programming 3 - Longest Alternating Subsequence}
\begin{multicols}{2}
\begin{definition}
    A sequence is alternating iff the subsequence alternates from increasing and decreasing for one term. 
\end{definition}

\begin{leftbar}
    \textbf{Longest Alternating Subsequence (LAS)} 

    Given a sequence, what is the length of its longest alternating subsequence. 
\end{leftbar}

    Note that the behavior of such subsequence depends on the values in the sequence. If the first element is small, it might be good to look for a larger value as the next element. Define now $\inc(i)$ and $\dec(i)$ as the length of the LAS which ends at index $i$, in increasing or decreasing direction, respectively. We do not wish to dictate the direction before we try out both, hence yielding the mutual recursion.  
    \begin{align*}
    \displaystyle
    \inc(i) &= 
    \begin{cases}
        1 & i = 1 \\
        1 + \max_{\substack{j < i \\ a_j < a_i}}\dec(j) & \textrm{otherwise}
    \end{cases} \\ 
    \dec(i) &= 
    \begin{cases}
        1 & i = 1 \\
        1 + \max_{\substack{j < i \\ a_j > a_i}} \inc(j) & \textrm{otherwise}
    \end{cases}
    \end{align*}
    This gives a formula for $LAS(i)$. 
    $$
    LAS(i) = \max_{1 \leq i \leq j}\{\max(\inc(i), \dec(i)\}.
    $$


\begin{leftbar}
    \textbf{The Knapsack Problem}

    Input:
    \begin{itemize}
        \item $n$ items, each with a value $v_i > 0$ and a weight $w_i > 0$, with $v_i, w_i \in \N$.
        \item Knapsack with a capacity of $W$. 
    \end{itemize}
    \textbf{Goal:} Pack knapsack for maximum value. 
\end{leftbar}
Note that this is similar to the WIS problem (from 2.1). Define $OPT(i,w)$ as the max value of items $1, \ldots, i$ with weight limit $w$. Therefore, the goal is to find $OPT(n, W)$. There are a few different cases. 
\begin{itemize}
    \item True optimal does not pick $i$, giving $OPT(i, w) = OPT(i - 1, w)$. 
    \item True optimal picks $i$, giving $OPT(i, w) = v_i + OPT(i -1, w - w_i)$.  
\end{itemize}
Note that we reduce the budget in the second case. 

$$
OPT(i, w) = 
\begin{cases}
    0 & i = 0 \\
    OPT(i - 1, w) & w_i > w \\
    \max\left(\substack{OPT(i - 1, w), \\ v_i + OPT(i - 1, w - w_i)}\right) & \textrm{else}
\end{cases}
$$
\begin{analysis}
    Here, the number of subproblems is $n \cdot W$, and the time per subproblem is only $O(1)$. This gives a total time of $O(nW)$. It is critical to note that even though $W$ is not an input size -- it's just a value -- it still shows up in the run time. 
\end{analysis}
\end{multicols}

\newpage
\subsection{Lecture 9: Dynamic Programming 4}

\begin{paracol}{2}
\begin{leftbar}
    \textbf{Maximum Sum Subarray}

    Given an array $A = [a_i]_{i \leq n}$ of $n$ integers, what is the maximum sum out of all subarrays? That is, find 
    $$
    \max_{i \leq j}\sum_{k = i}^j a_k
    $$
\end{leftbar}
We have seen a few different ways to approach this problem: 
\begin{itemize}
    \item \textbf{Naive algorithm} cubic time (3 nested loops)
    \item \textbf{Precomputed sums} quadratic time (innermost loop is replaced with recomputed sum)
    \item \textbf{Divide and conquer} takes $O(n\log n)$ time. 
\end{itemize}

Define $MMS(j)$ to be the maximum subarray that ends with the element $a_j$. At the position $j$, we have a choice. We either look at $MMS(j)$ and find it useful and use it after adding $a_j$, or we disregard it and simply use $a_j$. This works because either way, we include $a_j$ as an element. This yields the Bellman equation 
$$
MSS(j) = 
\begin{cases}
    a_j & j = 1 \\
    \max\{a_j, MSS(j - 1) + a_j\} & \textrm{otherwise}
\end{cases}
$$
To find the true $MSS$, we take $\max_{1 \leq j \leq n} \{MSS(j)\}$. 

The algorithm \ref{MSS} is also a space-efficient way to compute $MSS$. 

On the right we also present some space-efficient algorithms for previous DP problems.  

\switchcolumn
\begin{algorithm}
    \caption{Saving Memory with MSS}
    \label{MSS}
    \begin{algorithmic}[1] 
    \Require An array $A$ of $n$ integers. 
    \Ensure The maximum possible sum over all subarrays over $A$. 
    \State $curr = \max (0, a_1)$
    \State $best = -\infty$
    \For{$j \in [2, \ldots, n]$}
        \State $curr = \max(a_j, a_j + curr)$
        \State $best = \max(best, curr)$
    \EndFor
    \State \Return $best$
    \end{algorithmic}
\end{algorithm}
\bigskip

\begin{algorithm}
        \caption{Saving Memory with Fibonacci}
        \label{Fibonacci}
    \begin{algorithmic}[1]
    \Require An integer $n \geq 0$. 
    \Ensure The $n$th Fibonacci number $F_n$. 
    \State $a,b = 0,1$
    \State $f = null$
    \For{$i = 2, \ldots, n$}
        \State $f = a + b$
        \State $a, b = b, f$
    \EndFor
    \State \Return $f$. 
    \end{algorithmic}
\end{algorithm}
\bigskip
\begin{algorithm}
\caption{Saving Space with Max GPS}
    \begin{algorithmic}[1]
        \Require $n \times m$ grid $A$ of positive numbers. 
        \Ensure Maximum path sum $\sum_{(i, j) \in P}A[i, j]$ in a path $P$ from $A[1, 1]$ to $A[n, m]$. 
        \State Initialize two $n + 1$ length arrays $curr$ and $prev$ to all $0$. 
        \For{$i = 1, \ldots, n$}
            \For{$j = 1, \ldots, m$}
                \State $curr[j] = A[i, j] + \max\{prev(j), curr(j - 1)\}$
            \EndFor
            \State $prev = curr$
        \EndFor
        \State \Return $curr[m]$. 
    \end{algorithmic}
\end{algorithm}
\switchcolumn

Here we look at the DP recurrence of the DroDroDur. Define $M(f, d)$ to be the minimum drops needed to find floor $x$ if there are $f$ floors and $d$ drones still left. We would like to find the minimum number of drops required, as an actual number. We have 

\switchcolumn
\bigskip
Below is the DroDroDur recurrence, with the output of $M(n, k)$. 
$$
M(f, d) =
\begin{cases}
    0 & f = 0 \\
    f & d \leq 1 \\
    1 + \min_{1 \leq j \leq f} \Bigl\{\max{\left(\substack{M(j - 1, d - 1) \\ M(f - 1, d)}\right)}\Bigl\} & \textrm{otherwise}
\end{cases}
$$
\switchcolumn
\end{paracol}


\newpage
\subsection{Lecture 10: Greedy Algorithm}
\begin{multicols}{2}
    \begin{leftbar}
        \textbf{Unweighted Scheduling Problem}

        \begin{itemize}
            \item Input: $n$ intervals $(s(i), f(i))$ with weights $v_i = 1$
            \item A compatible schedule $S \subset \{1, \ldots, n\}$ with no $i, j \in S$ overlap. 
        \end{itemize}
    \end{leftbar}
    Note that this is similar to the weighted interval scheduling problem, except here, the weights are just $1$ for each job. The runtime for the DP algorithm we came up with for WIS was $O(n\log n)$ time, but we may not need such a complicated algorithm for this one, since all we are trying to do is to pick the most number of jobs. 

    All greedy algorithms have the following pattern 
    \bigskip
    \begin{algorithm}[H]
        \caption{Greedy Algorithm Pattern}
        \label{GreedyPattern}
        \begin{algorithmic}
        \State Sort choices by \textbf{some method}
        \State $S = \{\}$
        \For{$1 \leq i \leq n$}
            \If{$i$ is compatible}
                \State $S = S \cup \{i\}$
            \EndIf
        \EndFor
        \State\Return $S$
        \end{algorithmic}
    \end{algorithm}

    According to algorithm \ref{GreedyPattern}, we need to sort the jobs based on some criteria. There are a few possible ones: 
    \begin{itemize}
        \item Interval length
        \item Earliest start time
        \item Fewest conflicts
        \item \textbf{Earliest finish time}
    \end{itemize}
    The correct greedy rule here is that we sort by the earliest finish time, with the intuition that we want to free up resources as soon as possible. 

    The intuition needs to be tested for correctness. We here present a proof. 
    \begin{prop}
        To solve the Unweighted Scheduling Problem, we need to sort the jobs by their earliest finish time. 
    \end{prop}
    \begin{proof}
        This proof is called \say{Greedy stays ahead}, meaning that when the greedy $G$ completes $k$ intervals, no other scheduling can complete equal to or more than $k$ intervals. 
        
        Consider $G = \{i_1, \ldots, i_r\}$ to be the greedy schedule and $O = \{j_i,\ldots, j_r\}$ to be a supposedly another optimal schedule. We need to show that for all indices $t$, we have $f(i_t) \leq f(j_t)$, where $f$ is the finish time. Assume again for contradiction that there exists some interval $k$ that $O$ picks such that $f(i_r) \leq f(j_r) \leq s(k)$. If such $k$ exists, then Greedy is not optimal. However, such $k$ cannot exist, since Greedy searches over every single job and finds all compatible ones. So no such $k$ can exists. 

        Here, we just need to prove that $f(i_t) \leq f(j_t)$ for all $1 \leq t \leq r$, which can be done by induction. It is easy to see for the base case that $f(i_1) \leq f(j_1)$, since it is by the design of the algorithm that Greedy always chooses the earliest finish time. Now we assume that $f(i_t) \leq f(j_t)$ for all $t \leq k$ and try to prove $f(i_{k + 1}) \leq f(j_{k + 1})$. The next Greedy choice $i_{k + 1}$ is the earliest-finishing interval that starts after $f(i_k)$. Since $f(i_k) \leq f(j_k)$, we have $s(j_{k+1})$ is also after $f(i_k)$. Therefore, when Greedy picked $i_{k + 1}$, the option $j_{k + 1}$ was available. The Greedy chooses the one that finishes earliest, so $f(i_{k + 1}) \leq f(j_{k + 1})$. 
        
        
        Consider the following inequalities: 
        \begin{itemize}
            \item $f(i_k) \leq f(j_k) \leq s(j_{k + 1}) \leq f(j_{k + 1})$
            \item $f(i_k) \leq s(i_{k+1}) \leq f(i_{k + 1})$. 
        \end{itemize}
        
    \end{proof}
\end{multicols}

\newpage
\subsection{Lecture 11: Greedy Algorithm 2}
\begin{multicols}{2}
\begin{leftbar}
    \textbf{Minimum Lateness Scheduling}

    Given $n$ jobs with length $t_i$ and unique deadlines $d_i$. The lateness of job $i$ is defined as $L(i) = \max(0, f_i - d_i)$. We want to minimize the lateness of the schedule $S$ by minimizing $L(S) = \max_iL(i)$
\end{leftbar}
Note that technically we need to find an ordering of jobs, because we can start a job immediately after we finish the previous one. The optimal solution will not have any idle time. There are a few possible schedulings: 
\begin{itemize}
    \item Shortest lengths first
    \item \textbf{Earliest deadline first}
    \item Minimum slack $d_i - t_i$ first
\end{itemize}
In fact, earliest deadline first is the correct answer. 
\bigskip
\begin{algorithm}[H]
\caption{Minimum Lateness Scheduling}
    \begin{algorithmic}[1]
        \State Sort input so that $d_1\leq d_2 \leq \cdots \leq d_n$ 
        \State Initialize $s$, $f$ arrays for start and finish times
        \State $r = 0$
        \For{$1 \leq i \leq n$}
            \State $s[i] = r$
            \State $f[i] = s[i] + t_i$
            \State $r = f[i]$
        \EndFor
        \State\Return $s, f$
    \end{algorithmic}
\end{algorithm}

We present a proof of optimality, known as the exchange argument. 
\begin{proof}
    We give an outline of the proof first: Let $\mathcal{O}$ be an optimal schedule. We show that $L(\mathcal{O}) \geq L(G)$. Starting with $\mathcal{O}$, incrementally make changes to make it more like $G$. This gives a sequence of schedules $\mathcal{O} = \mathcal{O}_1, \ldots, \mathcal{O}_m = G$. At each step we will show that $L(\mathcal{O}_i) \geq L(\mathcal{O}_{i+1})$, which means that the changes can only improve the schedule. Then, a chain of inequalities will show that $L(\mathcal{O}) = L(\mathcal{O}_1) \geq \cdots\geq L(\mathcal{O}_m) = L(G)$. 

    Now we begin the proof. Define a pair of jobs $i, j$ such that $d_i < d_j$ but job $j$ is earlier in schedule than job $i$. We call this pair an \textbf{inversion}. Note that Greedy algorithm $G$ has no inversions by design. Suppose now each step from $\mathcal{O}_i$ to $\mathcal{O}_{i+1}$ is removing an inversion. There can be at most ${n\choose 2}$ inversions in the case that every pair of jobs is inverted. Therefore, within ${n\choose 2}$ steps, every $\mathcal{O}$ can be converted to some $G$. If we can prove that $L(\mathcal{O}_i) \geq L(\mathcal{O}_{i+1})$, in other words, removing inversions does not increase lateness, then we're done. 

    Proving that removing inversions does not increase lateness involves $3$ steps: 
    \begin{itemize}
        \item \textbf{If there is an inversion, then there is a consecutive inversion. }

        Suppose that there exists an inversion pair $i < j$ such that $d_i > d_j$. If $j = i + 1$, then we're done. Otherwise, we look at $d_{i + 1}$. If $d_i > d_{i + 1}$, then we're done. If not, we look at the pair $d_{i + 1}$ and $d_{i + 2}$. If $d_{i + 1} > d_{i + 2}$, we're done. If not, we keep on going. In the worst case, we get an increasing sequence of $d_j < d_i < d_{i + 1} < d_{i + 2} < \cdots < d_{j - 1}$, then we have found a consecutive inversion pair of $d_{j - 1}$ and $d_j$. 
        \item \textbf{Swapping inverted jobs reduces an inversion.}

        This is true by definition of inversion pairs. 
        \item \textbf{Swapping consecutive inversions does not increase lateness.} 

        Suppose that the pair $j, i$ was a consecutive inversion (in other words, $d_i < d_j$) with $s_j < f_j = s_i < f_i$. Now, we swap the inversion pair and get $s_j = s'_i < f'_i = s'_j < f'_j = f_i$. Since the order is $j, i$ before the jobs were swapped, after the swapping, the lateness of $i$ cannot increase since it is now being done earlier than before. Now, $l_j = f_j - d_j$ and $l'_j = f'_j - d_j$, but 
        \begin{align*}
            l'_j &= f'_j - d_j \\
            &= f_i - d_j \\
            &\leq f_i - d_i
        \end{align*}
        This shows that after swapping of consecutive inversion, the lateness of $i$, of $j$, and of $k$ for any $k \neq i, j$ does not increase.
    \end{itemize}
    This proves the inequality that 
    $$L(\mathcal{O}) = L(\mathcal{O}_1) \geq \cdots\geq L(\mathcal{O}_m) = L(G)$$
    where there are $0 \leq m \leq {n \choose 2}$ number of inversions. 
\end{proof}
Note that with $n$ jobs, the algorithm minimizes lateness in $O(n \log n)$. 

\begin{leftbar}
    \textbf{Fractional Knapsack Problem}

    We have $n$ items, each with a value of $v_i > 0$ and a weight $w_i > 0$. We want to maximize the value while packing the knapsack with capacity $W$. Now we are allowed to pick a fractional amount of each item. 
\end{leftbar}

In other words, we want a list $[f_1, \ldots, f_n]$ denoting the fraction of each item picked. The knapsack constraint is 
$$
\sum_{1 \leq i \leq n}f_iw_i \leq W. 
$$
Our goal is to find
$$
\max \sum_{1 \leq i \leq n}f_iv_i
$$
The Greedy algorithm is fairly simple. 
\bigskip
\begin{algorithm}[H]
\caption{Fractional Knapsack Problem}
\begin{algorithmic}[1]
    \State Sort items in decreasing order of $v_i/w_i$. This represents the unit weight. 
    \State $w = W$, $f = [0, \ldots, 0]$
    \For{$1 \leq i \leq n$}
        \If{$w_i < w$}
            \State $f[i] = 1$, $w = w - w_i$
        \Else 
            \State $f[i] = w/w_i$, $w = 0$
        \EndIf
        \If{$w \leq 0$}
        \State break
        \EndIf
    \EndFor
    \State \Return $f$
\end{algorithmic}
\end{algorithm}
Now we prove that Greedy $G$ is optimal. 
\begin{proof}
    Suppose $G = [f_1, \ldots, f_n]$ is Greedy and $\mathcal{O} = [r_1, \ldots, r_n]$ is an optimal algorithm. Assume that the items are ordered in decreasing order of $v_i/w_i$, which is their unit weight. 

    Find the first spot $i$ such that $f_i \neq r_i$. Greedy implies that $f_i > r_i$. This means that there exists $j > i$ with $r_j > f_j$. In other words, the optimal solution has to use up the weight somewhere else. Consider 
    $$
    \mathcal{O}' = \{r_1, \ldots, r_i + \epsilon, \ldots, r_j - (\epsilon w_i)/w_j, \ldots, r_n\}
    $$
    The value of $\mathcal{O}'$ is the value of $\mathcal{O}$ plus $\epsilon v_i - ((\epsilon w_i)/w_j)v_j$, which is greater than $0$. This is a contradition since $\mathcal{O}$ is optimal. 
\end{proof}
\end{multicols}

%\pagebreak
\section{Graph Algorithms and Advanced Data Structures}

\subsection{Lecture 12: Basics of Graphs}

\begin{multicols}{2}
\begin{definition}[Graph]
A pair $G = (V, E)$ is a graph, where $V = \{v_1, \ldots, v_n\}$ is the set of \textbf{vertices}, and $E = \{e_1, \ldots, e_m\}$ is the set of \textbf{edges}, where $e_k = \{v_i, v_j\}$ is an edge connecting two vertices. 

The number of vertices ($n$) is called the \textbf{order}, and the number of edges ($m$) is called the \textbf{size}. 
\end{definition}
Note that the definition of graph we gave was the most generic one. For example, if a graph were to have directionality over the edges, we could change the definition of the edge to an ordered pair where $e_k = (v_i, v_j)$. 
Here is an example of a graph. 
$$
\psmatrix[colsep=2cm,rowsep=1cm,mnode=circle]
u&v&z\\
x&y
\ncline{-}{1,2}{1,3}
\ncarc[arcangle=-30]{-}{1,1}{2,2}
\ncline{->}{2,1}{1,3}
\endpsmatrix
$$
In the graph above, we have $V = \{u, v, z, x, y\}$, and $E = \big\{\{v, z\}, \{u, y\}, (x, z)\big\}$. 
\begin{definition}[Simple graph]
    A \textbf{simple graph} is a graph that does not have any self cycles or multiple edges between the same two nodes. 
\end{definition}
\textbf{Unless otherwise specified, we will be working with simple graphs in this course.}
Below is an example of a non-simple graph. 
$$
\psmatrix[colsep=3cm,rowsep=1cm,mnode=circle]
[name = A] a 
&b
\ncarc[arcangle = 30]{-}{1,1}{1,2}
\ncarc[arcangle = -30]{-}{1,1}{1,2}
\nccircle[angleA = 90, armA = 0.2cm]{-}{A}{0.5cm}
\endpsmatrix
$$
\begin{definition}[Basics of undirected graphs]We present a few definitions for undirected graphs. 
\begin{itemize}
    \item Two vertices $u, v$ in graph $G$ are \textbf{neighbors} of each other if $e = \{u,v\} \in E$. 
    \item The number of edges incident with a vertices $v$ is called the \textbf{degree} of $v$, denoted by $\deg v$. 
    \item A graph $H$ is a \textbf{subgraph} of $G$ if $V(H) \subset V(G)$ and $E(H) \subset E(G)$. We write $H \subset G$. 
    \item A $u-v$ \textbf{walk} $W$ in $G$ is a sequence of vertices in $G$, beginning with $u$, ending with $v$, and such that consecutive vertices in the sequence are adjacent. The \textbf{length} of the walk $W$ is the number of edges in the sequence. 
    \item A $u-v$ walk is \textbf{open} if $u\neq v$. A walk that is not open is \textbf{closed}. 
    \item A $u-v$ \textbf{trail} is a $u-v$ walk in which no \textit{edge} is traversed more than once.
    \item A $u-v$ \textbf{path} is a $u-v$ walk in which no vertices are \textit{repeated}. 
    \item A $k$ \textbf{cycle} is a closed trail with no repeating vertices. 
    \item Two vertices $u$ and $v$ are \textbf{connected} if there exists a $u-v$ path. A graph $G$ is \textbf{connected} if every pair of vertices is connected. A graph that is not connected is \textbf{disconnected}. 
    \item A connected subgraph of $G$ that is not a proper subgraph of another other connected subgraph of $G$ is called a \textbf{component} of $G$. The number of components of $G$ is written as $k(G)$. Therefore, a graph $G$ is connected iff $k(G) = 1$. 
    \item The complement $\bar{G}$ of $G$ is a graph with vertex set $V(G)$ and edge set $E(\bar{G})$ such that $uv \in E(\bar{G})$ iff $uv \notin E(G)$. 
\end{itemize}
    
\end{definition}

\begin{definition}[Basics of directed graphs]
    We present a few definitions of directed graphs
    \begin{itemize}
        \item A directed graph is \textbf{strongly connected} if for all $u,v \in V$, there exists a $u-v$ path and $v-u$ path. 
        \item A directed graph is \textbf{weakly connected} if when we ignore the directionality of the graph, it is connected in the unidirectional sense. 
    \end{itemize}
\end{definition}
\begin{definition}[Classification of graphs] We present some common classes of graphs. 
\begin{itemize}
    \item A graph is \textbf{complete} if every two distinct vertices are adjacent. We denoted this by $K_n$ where $\abs{V} = n$ is the order of the graph. The size of $K_n$ is $\abs{E(K_n)} = {n\choose 2}$. 
\end{itemize}
\end{definition}
\begin{theorem}[First Theorem of Graph Theory]
If $G$ is a graph of size $m$, then 
$$
\sum_{v\in V(G)}\deg v = 2m
$$
    
\end{theorem}
\end{multicols}


\subsection{Lecture 13: I missed}

\newpage 

\subsection{Lecture 14: BTS, DFS, and Topological Ordering}
\begin{multicols}{1}
\begin{definition}
    A \textbf{binary search tree} (BTS) is a rooted tree data structure with the invariance: 
    \begin{itemize}
        \item The key of a node is greater than all keys in the left subtree.
        \item The key of a node is smaller than all keys in the right subtree. 
    \end{itemize}
\end{definition}
Here's an example of BTS
$$
\psmatrix[colsep=0.5cm,rowsep=0.5cm,mnode=circle]
  &   &   & 4 \\
  & 2 &   &   &   & 8 \\
1 &   & 3 &   & 7 &   & 9
\ncline{->}{1,4}{2,2}
\ncline{->}{1,4}{2,6}
\ncline{->}{2,2}{3,1}
\ncline{->}{2,2}{3,3}
\ncline{->}{2,6}{3,5}
\ncline{->}{2,6}{3,7}
\endpsmatrix
$$
We can do the following operations on BST: 
\begin{itemize}
    \item Search for a key $k$
    \item Insert new key $k$
    \item Delete existing key $k$
\end{itemize}
Every operation depends on time required for searching, which takes at best $O(\log n)$ time in the case that the tree if \say{balanced}. \textbf{Self-balancing} BTS's will re-balance itself after each operation, keeping each searching, inserting, and deleting in $O(\log n)$ time. 

Some applications of BST include
\begin{itemize}
    \item \textbf{Sets} good for searching for elements. 
    \item \textbf{Associative arrays} a set of keys associated with a value. 
    \item \textbf{Priority queues} The tree stories the values of priorities, allowing for quick access to highest priority values. 
\end{itemize}
Note that the $\texttt{clock}$ variable means how many edges we have explored. It is bounded by $2n$, where $n$ is the number of vertices. 
\begin{algorithm}[H]
\caption{Depth First Search}
    \begin{algorithmic}[1]
    \Require  
    \Ensure 
    \State $G = (V, E)$ is a diagraph
    \For{$u \in V$}
        \State $p[u], d[u], f[u] = \texttt{null}, -1, -1$.
    \EndFor
    \State $\texttt{clock} = 1$
    \Procedure{\texttt{DFS}}{$u$}
        \State $d[u] = \texttt{clock}$
        \State $\texttt{clock} ++$
        \For{$(u,v) \in E$}
            \If{$d[v] = -1$}
                \State $p[v]= u$
                \State $\Call{\texttt{DFS}}{v}$
            \EndIf
        \EndFor
        \State $f[u] = \texttt{clock}$
        \State $\texttt{clock}++$
    \EndProcedure
    \end{algorithmic}
\end{algorithm}

There are different kinds of edges that $\texttt{DFS}$ could return in the \textbf{tree}. These terms are relative and depends on the root we started. 
\begin{itemize}
    \item \textbf{Tree} edges: edges that discover new nodes. 
    \item \textbf{Forward} edges: ancestor to descendant.
    \item \textbf{Back} edges: descendant to ancestor.
    \item \textbf{Cross} edges: no ancestral relation between vertices. 
\end{itemize}
This ordering of edges are in descending order of priority, meaning that if the edge $(a,b)$ is a forward edge (with $a$ being the parent of $b$), yet we discovered $b$ through $a$, the edge $(a,b)$ is still labeled as a tree edge. 

For any directed edge $(u,v)$, we have 
\begin{itemize}
    \item $d[u] < d[v] < f[v] < f[u] \implies$ \textbf{tree} or \textbf{forward}
    \item $d[v] < d[u] < f[u] < f[v] \implies$ \textbf{back}
    \item $d[v] < f[v] < d[u] < f[u] \implies $ \textbf{cross}
\end{itemize}
The implementation of BST for directed and undirected graphs are the same. However, due to the lack of directionality, we only have tree of back edges in undirected graphs. (\textbf{Exercise}). 
\begin{definition}
    A directed acyclic graph (DAG) is a directed graph with no directed cycles. 
\end{definition}
Since DAG is acyclic, there is some node ordering such that all edges go left to right. Such ordering is called an \textbf{topological ordering}. 
\begin{definition}[Topological Ordering]
    Give a digraph $G = (V, E)$, a topological ordering is a sequence of vertices $v_1, \ldots, v_n$ such that for every directed edge $u \to v \in E$. 
    In other words, the node $u$ appears before $v$ in the ordering. 
\end{definition}

Note that only DAG's can have topological ordering. 

We have the following algorithm
\begin{itemize}
    \item First node has no incoming edges
    \item Repeat
    \begin{itemize}
        \item Find node $u$ with incoming degree $0$.
        \item Append $u$ to current TopOrdering.
        \item Update in-degrees
    \end{itemize}
\end{itemize}
\end{multicols}

\subsection{Lecture 14: TopOrders, SCCs, BFS, Shortest Paths}
\begin{multicols}{2}
    Another topological ordering implementation uses DFS. It relies on the fact that a DAG does not have back edges. We can order by decreasing DFS finish time. This can be done in linear time (\textbf{Exercise}). 
    \begin{definition}[Connected Component]
    A connected component is maximal subgraph $V' \subset V$ that is connected in $G = (V, E)$. 
    \end{definition}
    We can find all connected components of a graph by the following algorithm. 
    \begin{enumerate}
        \item Pick any new node $v$ and run a \texttt{DFS($v$)}. 
        \item Everything reached is a newly found connected component. 
        \item If any node $u$ is not in that component, repeat by \texttt{DFS($u$)}. 
    \end{enumerate}
    We can find all connected components in $\Theta(n + m)$ time. Now we introduce the concept of strongly connectedness in digraphs. 
    \begin{definition}[SCC]
        A strongly connected component of a digraph is the maximal subset $V'$ of vertices such that every vertex in $V'$ is reachable from every other vertex in the set. Formally, a set $V' \subset V$ is a SCC if for all $u, v \in V'$
        $$
        u \rightsquigarrow v \textrm{ and } v \rightsquigarrow u
        $$
        where $\rightsquigarrow$ means there is a directed path to. 
    \end{definition}

    \begin{theorem}[SCC Decomposition]
        Every directed graph can be compressed into a DAG of SCC's. 
        $$
        G \to G^{SCC} = (V^{SCC}, E^{SCC})
        $$
        where each $SCC$ is a node in $V^{SCC}$, and 
        $$
        E^{SCC} = \{(c_1, c_2) | (u,v) \in E \land u \in C_1, v\in C_2\},
        $$
        where $C_1, C_2 \in V^{SCC}$. 
    \end{theorem}
    Note that this digraph must be acyclic because if there was a cycle, it would have been merged into a component in SCC. This DAG is called the condensation graph and has no cycles. 
    \begin{definition}[Transpose]
        The transpose of the graph $G = (V, E)$ is denoted as $G^T = (V, E^T)$, where $E^T = \{(v,u) | (u, v) \in E\}$. This flips the direction of every edge. 
    \end{definition}
    Note that $G$ and $G^T$ have the same SCC, with directions flipped. 
    
    Now we transition to discussing BFS. Suppose we start with \texttt{BFS$(s)$} for $s \in V$. We have the following steps. 
    \begin{description}
        \item [Step 0] $L_0 = \{s\}$
        \item [Step 1] $L_1 = $ all neighbors of $L_0$
        \item [Step 2] $L_2 = $ all neighbors of $L_1$ that are not in $L_0$ or $L_1$. 
        \item [Step k] $L_k = $ all neighbors of $L_{k - 1}$ that are not in $L_0, \ldots, L_{k - 1}$. 
    \end{description}
    
    Below we present the algorithm for BFS. 
    \begin{algorithm}[H]
        \caption{\texttt{BFS($G = (V, E), s$)}}
        \begin{algorithmic}[1]
            \For{$v \in V$}
                \State $found[v] = false$
                \State $layer[v] = \infty$
            \EndFor
            \State $found[s] = true$ and $layer[s] = 0$
            \State Let $i = 0$, $L_0 = s$, and $T = \emptyset$
            \While{$L_i$ is not empty}
                \State Initialize new layer $L_{i + 1}$
                \For{$u \in L_i$}
                    \For{$(u,v) \in E$}
                        \If{$found[v] = false$}
                            \State $found[v] = true$
                            \State $layer[v] = i+1$
                            \State $T \leftarrow T\cup\{(u,v)\}$
                            \State $L_{i + 1} \leftarrow L_{i+1} \cup\{v\}$
                        \EndIf
                    \EndFor
                \EndFor
                \State $i \leftarrow i+1$
            \EndWhile
        \end{algorithmic}
    \end{algorithm}
    \begin{definition}[Distance]
        Distance $d(u,v)$ is the number of edges in shortest path from $u$ to $v$. Note that $M = (G, d)$ is a metric space for undirected $G$. 
    \end{definition}
    Often times, the notion of distance needs to take into account the weights of each edge. We would like to minimize or maximize it based on our needs. Therefore, the distance of a path can now be defined as a sum of edge weights. 
    $$
    d(u,v) = \sum_{e_i \in u \rightsquigarrow v}w(e_i)
    $$
    There are three different types of shortest path problems. 
    \begin{itemize}
        \item \textbf{Shortest path} Give $s,t \in V$, find the shortest path from $s$ to $t$. 
        \item \textbf{Single-source} Give $s \in V$, find shortest path from $s$ to every $t \in V$. 
        \item \textbf{All-pairs} Find shortest path between every pair in $V$. 
    \end{itemize}
    Sadly, BFS will not suffice for finding shortest distance with weights. 
\end{multicols}


\subsection{Lecture 15: Shortest Paths in Graphs}
\begin{multicols}{2}
    In general, when dealing with shortest paths in graphs, we often have a weight function $w: E \to \R$ where $w(e) = w(u,v)$, and the graph is defined as $G = (V, E, w)$.  
    \begin{definition}[Length and Distance of Weighted Path]
        For a path $v_1 \rightsquigarrow v_k = (v_1, v_2, \ldots, v_{k - 1}, v_k)$, we define the length to be 
        $$
        \ell(v_1 \rightsquigarrow v_k) = \sum_{1 \leq i \leq k - 1} w(v_i, v_{i + 1})
        $$
        We define the distance between $u, v$ as the shortest length path. That is, 
        $$
        d(u,v) = \min_{u\rightsquigarrow v}\ell(u\rightsquigarrow v)
        $$
    \end{definition}
    Recall the three different types of \say{shortest path}. We focus on single source shortest path (SSSP), which can be done by Dijkstra's algorithm for non-negative edge weights. Here we present the big picture of Dijkstra's algorithm. 
    \begin{leftbar}
        \textbf{Dijkstra's Algorithm}
        \begin{enumerate}
            \item Maintain a set of explored nodes $X$. Every node $u\in X$ has a known value of $d(s,u)$. 
            \item For every node $v \notin X$, track an upper bound for $d(s,v)$ with $d(s,v) \leq d(s,u) + w(u,v)$ for $(u,v) \in E$ and $u \in X$. 
            \item Add \say{closest} unexplored node into set $X$. 
            \item Repeat until every node explored. 
        \end{enumerate}
    \end{leftbar}
    Consider the following graph. 
    $$
\psmatrix[colsep=2cm,rowsep=1cm,mnode=circle, nodesep=2pt]
  & b & d \\
s \\
  & c & e
\ncline{->}{1,2}{1,3}^{2}
\ncline{->}{2,1}{1,2}^{10}
\ncline{->}{2,1}{3,2}_{3}
\ncline{->}{3,2}{1,3}^{8}
\ncline{->}{3,2}{3,3}_{2}
\ncarc[arcangle=20]{->}{1,2}{3,2}>{4}
\ncarc[arcangle=20]{->}{3,2}{1,2}<{1}
\ncarc[arcangle=20]{->}{1,3}{3,3}>{9}
\ncarc[arcangle=20]{->}{3,3}{1,3}<{7}
\endpsmatrix
$$
Dijkstra's algorithm produces the following table. We start from the node $s$. 
\bigskip
\begin{table}[H]
    \centering
    \begin{tabular}{c|c c c c c}
     & $s$ & $b$ & $c$ & $d$ & $e$ \\
     \hline\hline
     $p(\cdot)$ & $\emptyset$ & $b$ & $s$ & $b$ & $c$ \\
     $d_0(\cdot)$ & 0 & $\infty$ & $\infty$ & $\infty$ & $\infty$ \\
     $d_1(\cdot)$ & 0 & 10 & 3 & $\infty$ & $\infty$ \\
     $d_2(\cdot)$ & 0 & 7 & 3 & 11 & 5 \\
     $d_3(\cdot)$ & 0 & 7 & 3 & 11 & 5 \\
     $d_4(\cdot)$ & 0 & 7 & 3 & 9 & 5
    \end{tabular}
    \caption{Table for Dijkstra's Algorithm}
\end{table}
Note that DA breaks down for graphs with negative weights. We now prove DA. 
\begin{proof}
    Invariant: at any point in the algorithm, $d[u]$ is optimal for any $u \in X$ and $d[v]$ is best possible distance via paths using only vertices in $X$. We use an induction proof. 

    For the base case of $d[s] = 0$, the optimality is trivial by definition. Now, suppose that $X$ has $k$ vertices, and node $v$ is next to be added to $X$, we need to show that $d[v]$ is optional. Now, the only way (that we have not considered) for $v$ to have a shorter path than for some path $d[u]$ and $u \in X$ is that if there were some $a\notin X$ with $(a,v) \in E$ such that $d(s\rightsquigarrow u) + d(u, v) \leq d(s \rightsquigarrow a) + d(a,,v)$. However, this is impossible since we assumed that $v$ is the closest to $X$. Therefore, $d(u, a)$ for all $u \in X$ and $a\notin X$ must be that $d(u, v) \leq d(u, a)$. Since weights non-negative, the distance $d(s\rightsquigarrow a) + d(a, v)$ cannot be shorter. 
\end{proof}

Below we present the pseudo-code for DA. 
\bigskip
\begin{algorithm}[H]
    \caption{Dijkstra$(V, E, w, s)$}
    \begin{algorithmic}[1]
        \For{$v \in V$}
            \State $d[u] = \infty$
            \State $p[u] = null$
        \EndFor
        \State $Q = V$, and $d[s] = 0$
        \While{$Q \neq \emptyset$}
            \State $u = Q.pop(\arg\min_{w \in Q} d[w])$
            \For{$(u,v) \in E$}
                \If{$d[v] > d[u] + w(u,v)$}
                    \State $d[v] = d[u] + w(u,v)$
                    \State $p[v] = u$
                \EndIf
            \EndFor
        \EndWhile
        \State \Return $d,p$
    \end{algorithmic}
\end{algorithm}
The runtime depends on time taken to decrease key in priority queue $T_{dec}$ and on time taken to find minium in priority queue $T_{min}$. With BST based priority queues, both are $O(\log n)$. The total runtime is $O(\abs{E}\cdot T_{dec} + \abs{V} \cdot T_{min}) = O(\abs{E}\lg \abs{V})$. 

Since DA cannot handle edges with negative weights, we introduce Bellman-Ford Algorithm (BFA). We use DP to solve the problem. Each subproblem $OPT(v,j)$ is the length of the shortest path from $s\rightsquigarrow v$ using no more than $j$ edges. There are two possibilities. 
\begin{itemize}
    \item Use an existing path $OPT(v, j - 1)$
    \item Use some incoming edge. In other words, $OPT(u, j - 1) + w(u, v)$ for some $(u, v) \in E$. 
\end{itemize}
This gives the optimality equation. 
$$
OPT(i, j) = 
\begin{cases}
    0 & v = s \\
    \infty & v\neq s, j = 0 \\
    \min\left\{\substack{ OPT(v, j - 1) \\ \underset{(u,v) \in E}{\min}\left(OPT(u, j - 1) + w(u,v)\right)}\right\} & \textrm{otherwise}
\end{cases}
$$
Consider the following graph
    $$
\psmatrix[colsep=2cm,rowsep=1cm,mnode=circle, nodesep=2pt]
  & b & e \\
s \\
  & c & d
\ncline{->}{1,2}{1,3}^{2}
\ncline{->}{2,1}{1,2}^{-1}
\ncline{->}{2,1}{3,2}_{4}
\ncline{->}{3,3}{3,2}_{5}
\ncline{->}{1,2}{3,2}<{3}
\ncline{->}{1,3}{3,3}>{-3}
\ncarc[arcangle=15]{->}{1,2}{3,3}^{2}
\ncarc[arcangle=15]{->}{3,3}{1,2}_{2}
\endpsmatrix
$$

BFA gives the following table. 
\bigskip
\begin{table}[H]
    \centering
    \begin{tabular}{c|c c c c c}
        $OPT(v, j)$ & 0 & 1 & 2 & 3 & 4 \\
        \hline
        $s$ & 0 & 0 & 0 & 0 & 0 \\
        $b$ & $\infty$ & $-1$ & $-1$ & $-1$ & $-1$ \\
        $c$ & $\infty$ & $4$ & $2$ & 2 & 2 \\
        $d$ & $\infty$ & $\infty$ & 1 & $-2$ & $-2$ \\
        $e$ & $\infty$ & $\infty$ & 1 & 1 & 1
    \end{tabular}
    \caption{Table Produced by BFA}
    \label{tab:placeholder}
\end{table}

There are 3 different places where we could optimize the space and time usage. 
\begin{itemize}
    \item Once we have filled in $d[v,j]$, we don't need the value $d[v, j - 1]$ anymore. This could bring space usage from $O(\abs{V}^2)$ to $O(\abs{V})$. 
    \item We only need to check edge $(u, v)$ if $d[u, \cdot]$ changed in previous iteration. 
    \item We can also stop the outermost loop $1\leq i \leq n - 1$ if no $d[u]$ changes during an iteration. 
\end{itemize}

This gives space of $O(\abs{V})$ and time $O(\abs{E} \cdot \abs{V})$. 

Note that in the case that there is a cycle with a net negative value, BFA can forever walk through the cycle to make the net weight decrease indefinitely. The existence of such negative cycle can be detected by running BFA $n$ times. We check if there exists a vertex $v$ such that $d[v,n] < d[v,n - 1]$. If such $v$ exists, then there is a negative cycle on the shortest $s\rightsquigarrow v$ path.
\end{multicols}

\newpage
\subsection*{Lecture 16: Minimum Spanning Trees}

\begin{multicols}{2}

\begin{leftbar}
    Minimum Spanning Tree Problem

    We want to build a cheap, connected graph by choosing edges to build. 
    \begin{itemize}
        \item \textbf{Input} A weighted undirected graph
        \item \textbf{Output} A connected graph $G_{MST}= (V, T)$ containing all notes such that 
        \begin{itemize}
            \item If $u\rightsquigarrow v$ was in $G$, then $u \rightsquigarrow u$ must be in $G_{MST}$. 
            \item The total cost, which is the sum of weights of all edges, is minimized. 
        \end{itemize}
    \end{itemize}
\end{leftbar}


$$
\psmatrix[colsep=2cm,rowsep=0.75cm,mnode=circle, nodesep=2pt]
  & a \\
b &   & c \\
  & f \\
d &   & e \\
  & g
\ncline{-}{1,2}{2,1}^8
\ncline{-}{1,2}{2,3}^5
\ncline{-}{2,1}{2,3}^{10}
\ncline{-}{2,1}{4,1}<{18}
\ncline{-}{2,1}{3,2}^2
\ncline{-}{2,3}{3,2}^{3}
\ncline{-}{2,3}{4,3}>{16}
\ncline{-}{3,2}{4,3}>{30}
\ncline{-}{4,1}{3,2}^{12}
\ncline{-}{4,1}{5,2}_4
\ncline{-}{4,3}{5,2}_{26}
\ncline{-}{3,2}{5,2}>{14}
\endpsmatrix
$$

Since we want a connected graph containing all edges with minimum weight, it must be a tree. There are a few different algorithms that all compute the MST. 

\begin{description}
    \item[Prim] Start with some vertex $s$. At each step, add the cheapest edge that grows the connected component. 
    \item[Kruskal] Start with $T = \emptyset$. Add edges in ascending weight order unless they form a cycle. 
    \item[Reverse-Kruskal] Start with $T = E$. Delete edges in descending weight order unless it disconnects the graph. 
    \item[Boruvka] Start with $T = \emptyset$. In each rourd, add cheapest edge out of each connected component. 
\end{description}
\bigskip
\begin{algorithm}[H]
\caption{Prim$(V, E, w)$}
    \begin{algorithmic}[1]
        \For{$v \in V$} 
            \State $val[u] = \infty$, and $p[u] = null$
            \State $val[s] = 0$, and $Q = V$
        \EndFor
        \While{$Q \neq \emptyset$}
            \State $u = Q.pop(\arg\min_{v \in Q} val[v])$
            \For{$(u,v) \in E$}
                \If{$v \notin Q$ and $val[v] > w(u,v)$}
                    \State $val[v] = w(u,v)$
                    \State $p[v] = u$
                \EndIf
            \EndFor
        \EndWhile
        \State\Return $p$. 
    \end{algorithmic}
\end{algorithm}
The time complexity is $O(\abs{E} \lg \abs{V})$. 
\bigskip
\begin{algorithm}[H]
\caption{Kruskal}
\begin{algorithmic}[1]
    \State $T = \emptyset$
    \For{$v\in V$}
        \State $MakeSet(v)$
    \EndFor
    \State $E_{s} = $ sort $e$ in increasing order of $w(e)$
    \For{$(u, v) \in E_s$}
        \If{$FindSet(u) \neq FindSet(v)$}
            \State $T = T\cup \{(u,v)\}$
            \State $Union(FindSet(u), FindSet(v))$
        \EndIf
    \EndFor
    \State\Return $T$
\end{algorithmic}
\end{algorithm}

\begin{definition}[Cut]
    A cut $C$ is a subset of nodes $S\subset V$. A cut-set of a cut are edges that one endpoint in each cut. If a graph $G$ is cut into $C = (S, T)$, then the cut-set is $\{(u,v) \in E\ : u \in S, v\in T\}$. 
\end{definition}

\begin{theorem}[Cut and Cycle]
    The number of intersections between the cut-set and the cycles is always even. 
\end{theorem}

As a sketch of the proof, suppose there was a cut $C = (S,T)$. If $S$ or $T$ contains a cycle within itself, the number of intersection is $0$. If there was a cycle that starts from $S$, goes into $T$, and returns back to $S$, then it has to leave $S$ and come back, which contributes an intersection number of $2$. With this theorem comes come corollaries. 
\begin{itemize}
    \item If a cut-edge has minimum weight, we can \textbf{safely} pick it because it is optimal. 
    \item If an edge in a cycle has maximum weight, then it is \textbf{useless} since there are better alternative paths. This is because we can construct a cut from the cycle so that the maximum weight edge is in the cut set. By the previous point, we can pick an edge with a lower weight.  
\end{itemize}

Here we prove the correctness of Prim's algorithm. 
\begin{proof}
    Prim's algorithm always forms a cut set with the edges going out of the current minimum spanning tree. Since it chooses the lowest weight edge, it is always safe. 
\end{proof}
We now prove Kruskal's algorithm. 
\begin{proof}
    In Kruskal's algorithm, we either pick an edge in the cut set, or we pick an edge that's already in some component. However, Kruskal ignores all edges if forms a cycle when added, because if an edge forms a cycle, it must have a higher weight than other edges since they were already picked, making it useless. This means that we only add edges from the cut edge set. Note that the edge that we're adding is also safe, since by the definition of Kruskal, it is the edge with the smallest weight in all unseen edges, making it the safe edge. 
\end{proof}
\end{multicols}
%\pagebreak
\section{Flows, Linear Programming, and Complexity}

\subsection{Lecture 18: Introduction to Network Flows}

\begin{multicols}{2}
\begin{definition}
The flow network $N = (G, c(e), s,t)$ has directed $G = (V, E)$, edge capacities $c: E \to \R^{\geq 0}$, and the source $s$ and sink $t$. A flow is a path such that it uses some of the capacity along the edge. 
\end{definition}

A few common questions we wish to answer: 
\begin{itemize}
    \item Find flow through an edge $f : E \to \R^{\geq 0}$. How much can flow through? 
    \item Remove edges to stop flow from $s$ to $t$. What is the min capacity removed? 
\end{itemize}
Note a few things: 
\begin{itemize}
    \item The capacities are constraints. That is, $f(e) \leq c(e)$. 
    \item For all vertices except the source and sink, its in-flow must equal to the out-flow. That is, $$\sum_{e = (u, v)} f(e) = \sum_{e = (v,w)} f(e)$$
    \item The value of the flow is constant and defined as 
    $$
    \abs{f} = \sum_{e = (s,u)}f(e) - \sum_{e = (w,s)}f(e)
    $$
\end{itemize}

\begin{definition}
    A \textbf{cut} is a pair of set of vertices $(A, B)$ separates source $s \in A$ and sink $t \in B$. Note that $A \cup B = V$. The \textbf{cut capacity}, denoted by $\|A,B\|$, is the sum of capacities of edge from $A$ to $B$. 
\end{definition}

We now investigate more deeply. 
\begin{lemma}
    For any flow $f$ and any cut $(A,B)$, the value $\abs{f}$ equals net flow from $A$ to $B$. 
\end{lemma}
This lemma should intuitively make sense, since the total out-flow from the source should not disappear any time in the network. 
\begin{prop}
    We have a pair of dualities. 
    \begin{itemize}
        \item \textbf{Weak duality} Any flow value is upper bounded by any cut capacity. That is, 
        $$
        \abs{f} \leq \|A,B\|
        $$
        \item \textbf{Strong duality} Maximum flow is equal to the minimum cut capacity. That is, 
        $$
        \max_f \abs{f} = \min_{(A,B)} \|A,B\|
        $$
        \item This means that if $\abs{f} = \|A, B\|$, then $f$ is a max flow and $(A,B)$ is a min cut. 
    \end{itemize}
\end{prop}
The strong duality is true since $\abs{f}$ must be bounded above by every cut capacity, meaning that the maximum flow can be reaches is the minimum cut capacity. A few more useful notes: 
\begin{itemize}
    \item Maximum flows and minimum cuts can be computed in $O(\abs{V}\cdot \abs{E})$ time. 
    \item If all capacities are integer $c : E \to \Z^+$, then there is a max-flow with $f: E \to \Z^+$. 
    \item An edge is \textbf{saturated} if $f(e) = c(e)$. 
    \item An edge is \textbf{avoided} if $f(e) - 0$. 
    \item There's always a max flow in which only one of $f(u,v)$ and $f(v,u)$ is non-zero. 
\end{itemize}

In general, the term \textbf{reduction} means that we are converting a problem to another. Max flow and min cut are very good algorithms to reduce to, since they can help solving a large variety of optimization problems. 

\begin{leftbar}
    \textbf{Bipartite Matching}
    \begin{itemize}
        \item \textbf{Input} A bipartite graph $G = (V = L \cup R, E)$. 
        \item \textbf{Output} A subset of edges such that no vertex is matched with more than one edge. 
    \end{itemize}
\end{leftbar}
We can essentially reduce bipartite matching problem into a network flow problem. For vertices, we keep all old ones and add $s$ and $t$. Every edge $\{l, r\} \in E$ becomes a directed edge $(l,r) \in E$ with $\infty$ capacity. We added edges $(s,l)$ for all $l \in L$ and $(r, t)$ for all $r \in R$ with capacity $1$. Note that there is matching size $\abs{f}$ where $\abs{f}$ is the maximum flow value. Edges with $f(l,r) = 1$ give the actual matching edges where $l \in L$ and $r\in R$. 
\end{multicols}

% (Lessons 0–4) Probability distributions, Bayes' theorem, Bayesian networks, and the linear algebra basics (dot products, covariance, projections) everything else builds on.

% Unit 2: Regression 
% (Lessons 5–8) k-NN, OLS, cross-validation, bias-variance tradeoff, and Ridge regression — your first complete "build it, evaluate it, improve it" modeling cycle.

% Unit 3: Classification 
% (Lessons 9–15) Metrics, gradient descent, logistic regression, softmax, Lasso/Elastic Net, and SVMs — different mathematical routes to the same goal of drawing decision boundaries.

% Unit 4: Trees & Dimensionality Reduction 
% (Lessons 16–21) Decision trees, random forests, entropy/information theory, PCA, and LDA/QDA — a shift from linear boundaries to splitting rules and structural compression.

% Unit 5: Clustering & Applications 
% (Lessons 22–25) K-means, hierarchical clustering, KDE, ethics, and the closing research application — finding structure without labels, then tying it back to real use.


\end{document}
